\documentclass{beamer}
\usepackage{ctex}
\usetheme{SimplePlus}
\setbeamertemplate{footline}{}

\usepackage{amsmath,amssymb,mathtools}
\usepackage{tikz}
\usetikzlibrary{arrows.meta}
\usepackage{tikz-cd}

\usepackage{unicode-math}
%\setmathfont{STIX Two Math}
%\setmathfont{TeX Gyre Termes Math}
%\setmathfont{Libertinus Math}
%\setmathfont{Fira Math}

\usepackage[
  hybrid,                            % 允许在 Markdown 中混用 LaTeX 命令
  texMathDollars                     % 启用 $...$ 和 $$...$$ 数学公式
]{markdown}

\usepackage{graphicx} % Allows including images
\graphicspath{{./images/}}

\title{Motivation of Kripke And Beth Semantics}

\date{}

\begin{document}

\begin{frame}
    % Print the title page as the first slide
    \titlepage
\end{frame}

\begin{frame}[fragile]{直觉主义逻辑的背景}
\begin{markdown}
20世纪初，布劳威尔提出直觉主义数学，强调数学对象必须通过构造获得。海廷后来把直觉主义逻辑形式化，形成了现代直觉主义命题逻辑和谓词逻辑的基础。

此时的问题是：

> 如果不把命题简单地看作真假值，怎样给出一种严格的语义？
\end{markdown}
\end{frame}

\begin{frame}[fragile]{Kripke 语义的发展}
\begin{markdown}
1950年代，Saul Kripke 提出了以“可能世界”和“信息增长”为核心的语义。

Kripke 语义带来了几个重要成果：

- 为直觉主义逻辑提供了清晰的模型论解释；
- 说明了为什么排中律、双重否定消去等经典原则不能普遍使用；
- 为证明直觉主义逻辑的可靠性和完备性提供工具；
- 能够构造反例模型，证明某些公式不能在直觉主义逻辑中证明；
- 后来推广到模态逻辑、时态逻辑、动态逻辑和计算机科学中的程序语义。
\end{markdown}
\end{frame}

\begin{frame}[fragile]{Beth 语义的发展}
\begin{markdown}
Evert Willem Beth 在20世纪中期发展了以树、路径和证明展开为核心的语义方法。

Beth 语义具有更明显的证明论色彩：

- 它把证明看作逐步展开的过程；
- 分支表示证明中可能出现的不同发展；
- bar 表示所有可能的发展最终都会达到某种证明状态；
- 它帮助说明直觉主义逻辑中析取、存在量词和证明对象的构造性含义。

后来，Beth 语义与 Kripke 语义、拓扑语义、Heyting 代数语义之间建立了更深入的联系。
\end{markdown}
\end{frame}

\begin{frame}[fragile]{作用}
\begin{markdown}
\textbf{1. 解释直觉主义逻辑}

它们把“证明才是真”转化为严格的语义条件。例如：

- 蕴含不是简单的真假函数，而是对未来信息状态的稳定要求；
- 否定不是“假值”，而是“不存在任何未来证明”；
- 存在命题要求能提供见证；
- 析取要求能指出或最终确定哪一边成立。
\end{markdown}
\end{frame}

\begin{frame}[fragile]
\begin{markdown}
\textbf{2. 证明逻辑的可靠性与完备性}

语义学可以证明：
$$
\vdash A \quad\Longrightarrow\quad \models A
$$

即：可证明的公式在所有模型中都成立。

完备性则反过来说明：
$$
\models A \quad\Longrightarrow\quad \vdash A
$$

即：如果一个公式在所有模型中都成立，那么它确实可以在逻辑系统中证明。
\end{markdown}
\end{frame}

\begin{frame}[fragile]
\begin{markdown}
\textbf{3. 证明某些经典原则不能使用}

例如，构造一个 Kripke 模型，使得：
$$
w\not\models P\lor\neg P
$$

就可以说明排中律不是直觉主义逻辑的普遍定理。

类似地，还可以研究：
$$
\neg\neg P\to P
$$

为什么在经典逻辑中成立，却不一定在直觉主义逻辑中成立。
\end{markdown}
\end{frame}

\begin{frame}[fragile]{总结}
\begin{markdown}
可以这样记忆：

> **Kripke 语义把直觉主义真理理解为“在所有未来信息状态中保持成立”；Beth 语义则把它理解为“沿着所有可能的证明展开路径，最终都能达到相应的证明结果”。**

更简洁地说：

- **Kripke：状态语义、偏序、未来世界；**
- **Beth：过程语义、树结构、证明路径；**
- **共同点：都描述构造性真理和信息增长；**
- **区别：Kripke 更静态和模型论，Beth 更动态和证明论。**
\end{markdown}
\end{frame}

\end{document}